% SPDX-License-Identifier: Apache-2.0
\section{A Card in Its Own Language}
\label{sec:x-sd}

An SD card in SPI mode is a few hundred kilobytes a second on one
data line, which is what the SPI master of Section~\ref{sec:x-spi}
gives. The card's own protocol is faster: a command line that carries
48-bit commands out and 48- or 136-bit responses back, and one or
four data lines that carry blocks of 512 bytes either way, each line
with a CRC-16 of its own, at 25 or 50~MHz once the card is up. Every
command carries a CRC-7, and so does every response but one.

\code{Sd} is a host for that protocol behind AXI-Lite, and it does one
command at a time with at most one block. A program writes the
argument and then the command word, and the command word says what
follows, or that nothing does: eighty clocks with the line held high
are what a card wants before its first command, and one bit sends
those alone. Otherwise it says what follows the command: no response, a short one or a long one; a block to read
after it or a block to write; whether to wait for the card's busy
line afterwards; and whether to check the response's CRC, which
\code{ACMD41}'s answer does not carry. The host sends the command
with its CRC-7 computed as the bits leave, takes the response into
four words, checks its CRC-7 by feeding the received CRC into the
same register and looking for zero, moves the block through a buffer
of 128 words that the program reads or fills through one register,
and says in \code{status} how it went: a response that never came,
one whose CRC was wrong, a block that never came or a card that never
left busy, and a block whose CRC was wrong or that the card refused,
each a bit. The card clock is the system's divided; a board starts at
400~kHz and moves to 25~MHz once the card is up.

The other end of the wires in the example is \code{SdCard}, a model of
a card: it decodes the commands a host needs to bring a card up and
move blocks, answers them with their CRCs, sends blocks with a CRC-16
a line and checks the ones it is sent, answers with the CRC status
token and a spell of busy, and can be told to spoil one CRC, which is
how the host's checks are tested. It is stepped on the wires every
cycle, so what is on the lines is the protocol. The run brings the
card up on one line, reads a block, moves to four lines, writes a
block and reads it back.

A multi-block transfer is one command, \code{CMD18} or \code{CMD25},
then data-only transfers a block at a time, then \code{CMD12}. A card
streams a read back to back, and keeping up with that at 25~MHz
through a register is what DMA is for; the host here has one buffer,
and the engines of Section~\ref{sec:x-dma} are what would feed it.

\begin{figure}[htbp]
\centering
\input{sd_timing.tex}
\caption{The start of \code{CMD0}: the card clock, the command line
driven while the host holds it, the phase, the bit count and the
CRC-7 gathering as the bits leave.}
\label{fig:x-sd}
\end{figure}

\lstinputlisting[caption={\code{ex\_sd.rs}. Run with \code{bazel run //lib/examples:ex\_sd}.},label={lst:x-sd}]{lib/examples/ex_sd.rs}

\lstinputlisting[style=ascii,caption={What \code{ex\_sd} printed when the build ran it: the card's answers on the way up, the block read on one line, and the block written and read back on four.},label={lst:x-sd-out}]{out_sd.txt}
